\documentclass[a4paper,11pt]{article}
\usepackage[T2A]{fontenc}
\usepackage[utf8]{inputenc}
\usepackage[russian]{babel}
\usepackage{amsmath,amssymb,booktabs,tabularx,xcolor,tikz,fancyhdr,geometry}
\geometry{top=23mm,bottom=22mm,left=20mm,right=20mm}
\definecolor{accent}{HTML}{008080}
\definecolor{darkaccent}{HTML}{004D4D}
\usepackage[colorlinks=true,linkcolor=darkaccent]{hyperref}
\setlength{\parindent}{0pt}
\setlength{\parskip}{5pt}
\linespread{1.12}
\setlength{\headheight}{23pt}
\pagestyle{fancy}
\fancyhf{}
\fancyhead[L]{\small Теория вероятностей}
\fancyhead[R]{\small Екатерина Гнедкова}
\fancyfoot[C]{\thepage}
\renewcommand{\headrule}{\hbox to\headwidth{\color{accent}\leaders\hrule height 0.4pt\hfill}}
\newcommand{\heading}[1]{\par\medskip{\large\color{darkaccent}#1}\par\smallskip{\color{accent}\hrule}\medskip}
\newcommand{\rubric}[1]{\par\smallskip\textbf{#1}\enspace}
\newcommand{\err}{\textcolor{red!65!black}{\textbf{Ошибка:}}\enspace}
\newcommand{\idea}{\textcolor{darkaccent}{\textbf{Ключевая идея:}}\enspace}
\newcommand{\result}{\textcolor{darkaccent}{\textbf{Итог:}}\enspace}
\newcommand{\task}[2]{\clearpage\phantomsection\label{task:#1}\heading{Задание № #1. #2}}
\begin{document}
\begin{titlepage}
\centering
\vspace*{25mm}
{\Large\color{darkaccent} МАТЕМАТИКА\par}
\vspace{16mm}
{\huge Ревью домашней работы\par}
\vspace{10mm}
{\LARGE\color{accent} Деревья вероятностей\\[3pt]и последовательные испытания\par}
\vspace{18mm}
{\large Подготовка к ЕГЭ по профильной математике\par}
\vfill
{\large Екатерина Гнедкова\par}
\vspace{9mm}
27 ноября 2026 г.\par
\vspace{18mm}
{\small Лёвин Артём Александрович\\эксклюзивно для Екатерины Гнедковой\par}
\vspace{13mm}
\begin{tikzpicture}
\draw[accent,line width=.65pt] (0,0) .. controls (2,.85) and (4,-.85) .. (6,0) .. controls (8,.85) and (10,-.85) .. (12,0);
\end{tikzpicture}
\end{titlepage}
\setcounter{page}{1}
\heading{Сводная таблица ошибок}
В таблицу включены задания, в которых выявлены ошибки в ходе решения. Номера ведут к подробному разбору.\par
\medskip
\renewcommand{\arraystretch}{1.35}
\begin{tabularx}{\linewidth}{@{}c >{\raggedright\arraybackslash}p{35mm} >{\raggedright\arraybackslash}X >{\raggedright\arraybackslash}p{22mm} >{\raggedright\arraybackslash}p{23mm}@{}}
\toprule
№ & Тема & Суть ошибки & Ваш ответ & Верный ответ\\
\midrule
\hyperref[task:2]{2} & Сложение вероятностей & Неверный общий знаменатель дробей & $0{,}5$ & $\frac{13}{24}$\\
\hyperref[task:3]{3} & Дополнение события & Ошибка в вычислении разности & $0{,}0369$ & $0{,}0396$\\
\hyperref[task:8]{8} & Повторные испытания & Преждевременное округление; неточные промежуточные значения & $4$ & $4$\\
\bottomrule
\end{tabularx}\par
\vspace{10mm}
\heading{Основное направление работы}
Сохраняйте точные значения вероятностей до последнего действия. На дереве умножайте вероятности вдоль одного пути и складывайте вероятности несовместных подходящих путей. При поиске минимального числа испытаний обязательно проверяйте соседнее меньшее число.
\clearpage
\thispagestyle{empty}
\vspace*{\fill}
\begin{center}{\Huge\color{darkaccent} Разбор ошибок}\end{center}
\vspace*{\fill}
\task{2}{Прогулка к озеру}
\rubric{Условие.} Турист на первой развилке равновероятно выбирает левую или правую дорожку. Слева из четырёх равновероятных продолжений три ведут к озеру, справа из трёх продолжений к озеру ведёт одно. Остальные ведут к парковке. Найдите вероятность попасть к озеру.
\rubric{Ваше решение.} Дерево и выражение для вероятности построены верно:
\[\left(\frac12\cdot\frac14\right)\cdot3+\frac12\cdot\frac13=\frac38+\frac16.\]
Далее записано $\frac{6+3}{18}=\frac9{18}=0{,}5$. Ваш ответ: $0{,}5$.

\err Неверное приведение к общему знаменателю. Последний правильный шаг: $\frac38+\frac16$. Первый неверный переход: замена этой суммы дробью $\frac{6+3}{18}$.

\rubric{Почему это неверно.} При сложении дробей знаменатель должен быть общим кратным обоих исходных знаменателей. Число 18 не делится на 8. Наименьшее общее кратное 8 и 6 равно 24. При умножении знаменателя на дополнительный множитель умножается и числитель: значение дроби сохраняется.

\idea Два способа попасть к озеру несовместны: через левую и через правую дорожки. Вероятности этих способов складываются.

\rubric{Исправление и полное решение.} Вероятность левой дорожки равна $\frac12$, а условная вероятность озера после неё --- $\frac34$. Справа соответствующие вероятности равны $\frac12$ и $\frac13$. Следовательно,
\[P(\text{озеро})=\frac12\cdot\frac34+\frac12\cdot\frac13=\frac38+\frac16.\]
Первую дробь домножаем на 3, вторую на 4:
\[\frac38+\frac16=\frac{3\cdot3}{8\cdot3}+\frac{1\cdot4}{6\cdot4}=\frac9{24}+\frac4{24}=\frac{13}{24}.\]
\rubric{Проверка.} Вероятность парковки:
\[\frac12\cdot\frac14+\frac12\cdot\frac23=\frac18+\frac13=\frac{11}{24},\qquad \frac{13}{24}+\frac{11}{24}=1.\]
\result $\boxed{\frac{13}{24}}$. Дерево можно сохранить; исправить требуется сложение дробей.

\rubric{Задача на закрепление.} Турист равновероятно выбирает одну из двух троп. На первой тропе из пяти равновероятных продолжений три ведут к озеру, на второй из четырёх --- одно. Найдите вероятность попасть к озеру.
\task{3}{Повторное исследование}
\rubric{Условие.} В образце опасного вещества нет. Исследование даёт ложный положительный результат с вероятностью $0{,}02$, отрицательный --- с вероятностью $0{,}98$. При первом отрицательном результате проводят второе исследование; при положительном прекращают проверку. Результаты независимы. Найдите вероятность хотя бы одного положительного результата.

\rubric{Ваше решение.} Записано $1-0{,}98\cdot0{,}98=0{,}0369$. Ваш ответ: $0{,}0369$.

\err Арифметическая ошибка в последнем равенстве. Выражение $1-0{,}98\cdot0{,}98$ составлено правильно. По записи невозможно установить, возникла ошибка при умножении или при вычитании; неверным является итоговое равенство.

\idea Событие «хотя бы один положительный результат» противоположно событию «оба результата отрицательные». Поэтому сумма их вероятностей равна единице.

\rubric{Исправление и полное решение.} Чтобы ни одного положительного результата не было, сначала должен появиться отрицательный результат, затем ещё один отрицательный. Вероятность этого пути равна произведению:
\[P(--)=0{,}98\cdot0{,}98=\frac{98\cdot98}{10000}=\frac{9604}{10000}=0{,}9604.\]
При вычитании удобно записать единицу с четырьмя знаками после запятой:
\[P(\text{хотя бы один }+)=1{,}0000-0{,}9604=0{,}0396.\]
\rubric{Проверка другим способом.} Положительный результат может появиться сразу либо после первого отрицательного. Эти пути несовместны. Складываем их вероятности:
\[P(+)+P(-+)=0{,}02+0{,}98\cdot0{,}02=0{,}02+0{,}0196=0{,}0396.\]
Оба метода приводят к одному значению. Кроме того, вероятность двух отрицательных результатов вместе с искомой вероятностью даёт единицу.

\result $\boxed{0{,}0396}$. Выбор противоположного события верен; требуется точное вычисление.

\rubric{Задача на закрепление.} При проверке исправного прибора вероятность ложного сигнала равна $0{,}03$. Если первая проверка не дала сигнала, прибор проверяют ещё раз; при сигнале проверку прекращают. Результаты независимы. Найдите вероятность хотя бы одного ложного сигнала.
\task{8}{Автоматическая система стрельбы}
\rubric{Условие.} Вероятность уничтожить ещё не поражённую цель при каждом выстреле равна $0{,}65$. Результаты независимы, после первого успеха стрельба прекращается. Найдите наименьшее число разрешённых выстрелов, при котором вероятность уничтожения цели составляет не менее $0{,}98$.

\rubric{Ваше решение.} Вероятности успеха на разных путях складываются. Для двух, трёх и четырёх выстрелов записаны соответственно $0{,}878$, $0{,}958$ и $0{,}986$. Ваш ответ: $4$.

\err Преждевременное округление. Последний точный шаг --- выражение $0{,}65+0{,}35\cdot0{,}65$. Первое неточное равенство: $0{,}65+0{,}35\cdot0{,}65=0{,}878$; точное значение $0{,}8775$. При дальнейшем сложении использовано уже округлённое число. Итоговое количество выстрелов получилось верным, однако вероятности искажены.

\idea При разрешённых $n$ выстрелах цель останется непоражённой только при $n$ последовательных промахах. Вероятность промаха равна $1-0{,}65=0{,}35$. Поэтому $P_n=1-0{,}35^n$.

\rubric{Исправление и полное решение.} Вычисляем без промежуточного округления:
\begin{align*}
P_1&=1-0{,}35=0{,}65,\\
P_2&=1-0{,}1225=0{,}8775,\\
P_3&=1-0{,}042875=0{,}957125,\\
P_4&=1-0{,}01500625=0{,}98499375.
\end{align*}
Для трёх выстрелов вероятность меньше $0{,}98$, для четырёх --- больше. При увеличении $n$ величина $0{,}35^n$ убывает, поэтому $P_n$ возрастает; все меньшие числа выстрелов также не подходят.

\rubric{Проверка исходным методом.} При трёх выстрелах сумма вероятностей успеха на первом, втором или третьем выстреле равна
\[0{,}65+0{,}35\cdot0{,}65+0{,}35^2\cdot0{,}65=0{,}957125.\]
При четвёртом добавляется $0{,}35^3\cdot0{,}65=0{,}02786875$, что даёт $0{,}98499375$. Результаты совпали.

\result $\boxed{4}$. При проверке порога сохраняйте точные значения и обязательно сравнивайте соседние количества испытаний.

\rubric{Задача на закрепление.} Независимая попытка подключения к серверу успешна с вероятностью $0{,}7$. После успеха попытки прекращаются. Какое наименьшее число попыток нужно разрешить, чтобы вероятность хотя бы одного успешного подключения была не менее $0{,}99$?
\end{document}
